% q0060.tex — Cartel Demolition
\question{
  id        = {0060},
  title     = {Cartel Demolition},
  source    = {OBECON},
  year      = {2026},
  round     = {Seletiva IEO -- Phase C},
  language  = {English},
  topic     = {Micro},
  subtopic  = {Oligopoly / Cournot / Price Ceiling},
  type      = {objective},
  statement = {%
    Obeconomia's construction industry consists of five identical Cournot
    competitors. Market demand is $P = 100 - Q$ and marginal cost is $MC = 20$
    for all firms. The government is considering imposing a price ceiling.
    Which price ceiling maximises total surplus?
  },
  choices   = {%
    \item \$0.
    \item \$20.
    \item \$40.
    \item \$60.
  },
  answer    = {B},
  solution  = {%
    Total surplus is maximised at the competitive equilibrium, where price
    equals marginal cost: $P^* = MC = \$20$, yielding $Q^* = 80$ units.
    A price ceiling set at \$20 forces firms to price competitively, eliminating
    oligopoly markup and deadweight loss, thereby achieving the maximum possible
    total surplus. (B) is correct.
  },
}
